\documentclass[11pt]{article}
\usepackage[a4paper,margin=25mm]{geometry}
\usepackage{amsmath,amssymb,amsthm,booktabs}
\newtheorem{proposition}{Proposition}
\begin{document}
\title{Truncations of the Higman--Sims geometries}
\author{L091 classification component}\date{}\maketitle
We use the complete enumeration and the separate proof of bare incidence
isomorphism classes for the specified Higman--Sims action. There are
$1,6,30,41,8,1$ classes in ranks one through six and no higher ranks.
A truncation below means deletion of types from another geometry in this
classified family. The internal representative numbers refer to the
fixed file \texttt{full\_fusion03/result.g}.

\begin{proposition}
Exactly two rank two classes and ten rank three classes are truncations
of higher rank geometries in the family. No class of ranks four, five
or six is such a truncation. The sole rank one class is a truncation.
Consequently the numbers of classes which are not truncations are
\[
                     0,\quad4,\quad20,\quad41,\quad8,\quad1.
\]
\end{proposition}
\begin{proof}
Let $G$ have order $44\,352\,000$, and let $(G_i)_{i\in I}$ be a
base chamber's parabolic stabilizers. For a subset $J\subseteq I$,
the truncation is the coset geometry with parabolics $(G_i)_{i\in J}$.
It inherits flag transitivity. Thus its type sizes, pair-flag counts
and chamber count determine respectively the multisets
\[
\{|G_i|:i\in J\},\qquad
\{|G_i\cap G_j|:\{i,j\}\subseteq J\},\qquad
\left|\bigcap_{i\in J}G_i\right|.
\]
These are necessary invariants for incidence isomorphism with any
member of the classified list, since the group order is fixed.
Their use does not assume that an incidence isomorphism normalizes $G$.

For every one of the $86$ representatives of rank at least two, take
every proper subset $J$ with $|J|\geq2$. There are exactly
\[
 30\binom32+41\left(\binom42+\binom43\right)
 +8\left(\binom52+\binom53+\binom54\right)
 +\binom62+\binom63+\binom64+\binom65=756
\]
such subsets. For each, compare the displayed invariant with every
representative. Exactly fifteen comparisons survive. Each has a
saved permutation in the acquired full automorphism action of $G$
and a permutation of types, sending the actual truncation parabolics
onto those of its target. These give explicit coset isomorphisms.
There are no surviving comparisons without a positive witness.
The distinct targets are $3,4$ in rank two and
$8,13,18,20,21,22,28,29,30,31$ in rank three.

An arbitrary classified parent is incidence-isomorphic to a listed
representative; any such isomorphism transports its type subsets.
Hence the finite list of comparisons covers every possible parent
and truncation. Every excluded comparison fails an intrinsic incidence
invariant, while every surviving one is positively identified. This
proves completeness of the claimed twelve targets, including the
negative assertions in ranks four through six. The separately
checked rank one component supplies a truncation witness for its
unique $176$-object geometry. Subtraction from the complete rank
counts proves the displayed result.
\end{proof}

\paragraph{Attribution and computational scope.}
These are the truncation totals stated in D.~Leemans, \emph{The residually
weakly primitive geometries of HS}, Australasian Journal of Combinatorics
33 (2005), 231--236. Completeness of the parent classification is an
explicit acquired premise. The new finite comparison uses actual
subgroup intersections and the acquired simultaneous conjugation
method in GAP; every positive witness is checked directly. No inference
uses failure of a conjugator search. This component supplies the
truncation clause, not the supplementary diagrams, automorphism groups
or the two geometric constructions, and claims no new open result.
\end{document}
